%% English CV — Vo Ngoc Kha Ai. Compile: ./build.sh cv-en.tex
%% Mirrors cv-vi.tex; [square brackets] mark details still to be filled in.
\documentclass{cvchem}

\cvaccent{1F4E79}

% Squeeze onto one page if needed (1.0 default, 0.6 very tight):
\cvdensity{0.9}

\cvname{Vo Ngoc}{Kha Ai}
\cvtitle{Chemical Engineering Graduate --- QA/QC \& R\&D}
\cvquote{Final-year Chemical Technology student at IUH --- interned at a paint plant, hands-on with coating film test instruments in a QA/QC workflow.}

\cvcontact{\icmail}{khaai.iuh@gmail.com}{mailto:khaai.iuh@gmail.com}
\cvcontact{\icphone}{0822 909 029}{tel:+84822909029}
\cvcontact{\icpin}{Ho Chi Minh City, Vietnam}{}

\cvfooter{Vo Ngoc Kha Ai · Curriculum Vitae · 2026}

\begin{document}
\makecvheader

\section[\icmol]{Profile}
\cvsummary{%
Final-year \textbf{Chemical Engineering Technology} student at Industrial University
of Ho Chi Minh City (IUH), with a focus on \textbf{organic chemistry}.
Completed an internship at \textbf{Son Hai Van Paint Co., Ltd.}, covering the full
paint production chain --- raw materials (alkyd, acrylic, epoxy and polyester resins;
solvents; pigments and extenders), paste dispersion, fine grinding, let-down, and
final quality control. Operated coating film test instruments (viscosity, hiding power,
fineness, density, gloss, adhesion) and read results against standardised inspection forms. Thesis work focuses on \textbf{leather-like
composite materials from water hyacinth and agricultural by-products}. Able to use standard measurement and
basic analytical instruments in a university laboratory. Careful and accountable with
test data, and used to the rhythm of shift work in a plant. Seeking a
\textbf{QA/QC or R\&D} position in paint, coatings, materials or industrial chemicals.}

\section[\icgear]{Experience}
\cventry{Intern --- Technical Department (QA/QC)}{Son Hai Van Paint Co., Ltd. --- Supervisor: Le Quoc Thoi, Technical Manager}{Apr -- Jun 2026}{Tay Ninh}
\begin{cvitems}
  \item \textbf{Standards-based quality control}: carried out sampling and inspection across the \textbf{IQC --- IPQC --- FQC} stages on \textbf{400--500 kg} industrial batches; recorded results on standardised forms under an \textbf{ISO 9001:2015} system.
  \item \textbf{Operated coating film test instruments}: Stormer viscometer (KU), Pfund Cryptometer hiding power, Hegman fineness gauge, density cup, gloss meter (20°, 60°, 85°) and adhesion tester --- measured at \textbf{25\,±\,0.5\,°C} since KU viscosity is temperature-sensitive, and judged pass/fail against technical tolerances (e.g. density \textbf{1.25\,±\,0.02\,g/cm³}, alkyd primer \textbf{\textasciitilde70 KU}).
  \item \textbf{Familiar with long-duration durability testing} in the plant laboratory: \textbf{QUV} accelerated weathering, cylindrical mandrel bend (TCVN 2099/ISO 1519), impact resistance (TCVN 2100-2/ISO 6272-2) and \textbf{5\% NaCl salt spray at 35\,°C for 1000 h} (ISO 7253) --- including specimen preparation and rating of rust and blistering (ISO 4628-2) and scribe creep (ASTM D1654).
  \item \textbf{Colour and appearance}: used a colour matching cabinet and spectrophotometer to compare batches against customer standards --- \textbf{good colour discrimination, not colour-blind}.
  \item \textbf{Batch adjustment to spec}: understood the practice of adding solvent stepwise to bring viscosity onto target, never below spec, since an over-thinned batch is hard to recover.
  \item \textbf{Production line familiarity}: high-speed disperser, basket mill and hydraulic let-down tank --- understood how temperature, grinding time and resin feed rate affect batch quality.
  \item \textbf{Observed batch troubleshooting}: followed how plant engineers trace root causes of out-of-spec batches (resin shock from fast pumping, pigment flocculation, off-target density/viscosity/fineness) and adjust to recover a batch rather than scrap it.
  \item \textbf{Health, safety and environment}: followed fire-safety rules and PPE requirements, consulted MSDS, and segregated waste solvent and hazardous solid waste per plant procedure.
\end{cvitems}

\section[\icbook]{Education}
\cventry{B.Eng. in Chemical Engineering Technology}{Industrial University of Ho Chi Minh City (IUH) --- Faculty of Chemical Engineering}{2022 -- 2026}{Ho Chi Minh City}
\begin{cvitems}
  \item GPA \textbf{2.8}/4.0.
  \item \textbf{Graduation thesis}: \textit{Factors affecting the production of composite and leather-like materials from water hyacinth, agricultural by-products and waste paper from garment embroidery}. Supervisor: Dr. Pham Thi Hong Phuong.
  \item Core coursework: Organic Chemistry, Physical Chemistry, Reaction Engineering, Unit Operations, Instrumental Analysis, Process Safety.
\end{cvitems}

\section[\icflask]{Technical Skills}
\begin{cvfields}[.25\linewidth]
  \cvfield{Operated}{Stormer viscometer, Pfund Cryptometer, Hegman gauge, density cup, gloss meter, adhesion tester}
  \cvfield{Exposure to}{Salt spray chamber (ISO 7253), QUV weathering, Elcometer mandrel bend, impact tester; rust/blister rating ISO 4628-2, ASTM D1654}
  \cvfield{Lab instruments}{Able to use standard measurement and basic analytical instruments in a university laboratory}
  \cvfield{Materials}{Alkyd, acrylic, epoxy, polyurethane and polyester resins; solvent systems; pigments and extenders; natural-fibre composites (water hyacinth, agricultural by-products)}
  \cvfield{Quality systems}{ISO 9001:2015, 5S, IQC/IPQC/FQC workflow, SOPs, navigating TCVN --- ASTM --- ISO}
  \cvfield{HSE}{MSDS/SDS, PPE, fire safety, hazardous waste segregation and handling}
  \cvfield{Software}{ChemDraw, Excel (test data handling), Word --- technical reporting}
\end{cvfields}

\vskip.5ex
\begin{cvfields}[.55\linewidth]
  \cvmeter{Coating film testing}{4} & \cvmeter{Applying standards}{4}\\[.5ex]
  \cvmeter{Chemical safety (HSE)}{4} & \cvmeter{Lab instrumentation}{4}\\[.5ex]
  \cvmeter{Organic chemistry, materials}{3} & \cvmeter{Data handling, records}{3}\\
\end{cvfields}

\section[\iclab]{Research \& Projects}
\cventry{Thesis: Leather-like composites from water hyacinth and agricultural by-products}{Faculty of Chemical Engineering, IUH --- Supervisor: Dr. Pham Thi Hong Phuong}{2026}{Ho Chi Minh City}
\begin{cvitems}
  \item Investigating the process factors governing leather-like material made from \textbf{water hyacinth, agricultural by-products and waste paper} from garment embroidery.
  \item Aimed at valorising waste streams into an environmentally friendlier substitute for synthetic leather.
\end{cvitems}

\cventry{Industrial internship report (team of 8)}{Faculty of Chemical Engineering, IUH --- Supervisor: Le Nhat Thong, M.Sc.}{2026}{Ho Chi Minh City}
\begin{cvitems}
  \item Documented the complete paint manufacturing process: raw materials, block flow diagram, equipment, common failures and corrective measures.
  \item Compiled \textbf{12 quality tests} with their principles, apparatus, governing standards and step-by-step procedures.
\end{cvitems}

\section[\icaward]{Awards}
\begin{cvitems}
  \item \textbf{Second and Third Prize --- Provincial Science \& Engineering Fair} (social media category), high school.
  \item \textbf{Second Prize --- Provincial Chess Championship}, high school.
  \item \textbf{National Consolation Prize} --- \textit{Safe Traffic for a Smiling Tomorrow} competition.
  \item \textbf{Provincial Consolation Prize} --- Reading Culture competition.
\end{cvitems}

\section[\icweb]{Certifications}
\begin{cvfields}[.25\linewidth]
  \cvfield{IT}{Certificate in Basic Information Technology Application --- Industrial University of Ho Chi Minh City}
\end{cvfields}

\section[\icgear]{Other Activities}
\cventry{Volunteer Collaborator}{Ward 5 People's Committee, [City/Province]}{[dates]}{[location]}
\begin{cvitems}
  \item Supported local community programmes and events; teamwork and direct engagement with residents.
\end{cvitems}


\end{document}
